% Created 2026-06-12 Fri 20:25
% Intended LaTeX compiler: pdflatex
\documentclass{article}
\usepackage[utf8]{inputenc}
\usepackage[T1]{fontenc}
\usepackage{graphicx}
\usepackage{longtable}
\usepackage{wrapfig}
\usepackage{rotating}
\usepackage[normalem]{ulem}
\usepackage{amsmath}
\usepackage{amssymb}
\usepackage{capt-of}
\usepackage{hyperref}

\usepackage{fancyhdr}
\usepackage[top=25mm, left=25mm, right=25mm]{geometry}
\usepackage{longtable}
\fancyhead[R]{}
\setlength\headheight{43.0pt}
\usepackage[T1]{fontenc}
\usepackage[utf8]{inputenc}
\usepackage[spanish]{babel}
\usepackage[backend=biber, style=apa]{biblatex}
\addbibresource{/home/ani_quillu/ArquitecturaComputadores/FormatoTareas/bibliography.bib}
\author{Auz Juan Diego, Lanchimba Danny, Palomeque Matías, Quillupangui Ana María}
\date{\today}
\title{TallerIAS}
\hypersetup{
 pdfauthor={Auz Juan Diego, Lanchimba Danny, Palomeque Matías, Quillupangui Ana María},
 pdftitle={TallerIAS},
 pdfkeywords={},
 pdfsubject={},
 pdfcreator={Emacs 27.1 (Org mode 9.3)}, 
 pdflang={Spanish}}
\begin{document}

\maketitle
\tableofcontents

\fancyhead[C]{\includegraphics[scale=0.05]{../images/logoEPN.jpg}\\
ESCUELA POLITÉCNICA NACIONAL\\FACULTAD DE INGENIERÍA DE SISTEMAS\\
ARQUITECTURA DE COMPUTADORES}
\thispagestyle{fancy}


\section{Recursos}
\label{sec:org455adbc}
Para este taller consulte el libro \citetitle{stallings2006}
disponible en \href{https://epnecuador-my.sharepoint.com/:f:/g/personal/lenin\_falconi\_epn\_edu\_ec/EgjH2RoedD5NuswqpOt8ExsB\_DE052v9Rlrg0QpEtbimDg?e=WoqexR}{EPN-SHAREPOINT(Clic Me)} y use su WSL con Emacs para editar el archivo.
\subsection{Observaciones y Recomendaciones}
\label{sec:orgb0aee7d}
\begin{itemize}
\item Tenga en cuenta la ruta a las imágenes y a la bibliografía. Se le
\end{itemize}
recomienda que realice una clonación o un fork del repositorio de la
clase a fin de que no tenga problemas generando el documento.
\begin{itemize}
\item Instale \LaTeX completo en el computador que generará el archivo pdf.
\end{itemize}
\section{Instrucciones}
\label{sec:orgd4e8ff5}
Complete las solicitudes del taller y suba el archivo \texttt{.ORG} y \texttt{.PDF}
al aula virtual. Cuando trabaje en Emacs trate de utilizar las
combinaciones de teclas para navegar por el documento.

\section{Taller Máquina de Von Neuman}
\label{sec:orgb082f05}

La máquina de Von Neuman realiza operaciones en un ciclo repetitivo de
captación y ejecución. Para esto, el computador IAS usa una memoria de
40 bits. Cuando la memoria se usa numéricamente, se representa el
número en complemento a 2, reservando el bit inicial para el
signo. Cuando la memoria se usa para el registro de instrucciones, se
divide en dos partes de 20 bits. Los 8 primeros bits de cada parte
corresponden al \emph{opcode} y los restantes 12 bits de cada parte a la
dirección (operando) desde la que se debe leer o a la que se debe
escribir.

\subsection{Conjunto de Instrucciones IAS}
\label{sec:orgc72c80b}
Complete la Tabla 2.1 de instrucciones del computador IAS
\autocite[p.47]{stallings2006} usando la notación RTL\footnote{Register Transfer Language o Lenguaje de Transferencia de Registros}. Se suministra
la Tabla en código \LaTeX para tal efecto. Edite el archivo en Emacs
usando el modo principal ORG y cambie al modo \LaTeX para editar la
Tabla. El cambio de modo se realiza usando \texttt{M-x Latex-mode RET} y para
retornar al modo \texttt{ORG} haga \texttt{M-x org-mode RET}

\begin{itemize}
\item \texttt{M} equivale a \emph{Alt}
\item \texttt{x} es la tecla \emph{x}
\item \texttt{RET} es presionar \emph{enter}
\end{itemize}

Una vez activado el modo \LaTeX el siguiente código recibe colores
sobre las palabras clave.

\begin{table}[htbp]
  \caption{Instrucciones Maquina IAS}
  \label{tab:tablaISA-IAS}
  \centering
  \begin{tabular}{|cccc|}
    \hline
    Opcode   & Opcode Hex &  Simbolo          & RTL \\ \hline
    00001010 & 0x0A       &  LOAD MQ          & $[AC]\leftarrow[MQ]$\\
    00001001 & 0x09       &  LOAD MQ, M(X)    & $[MQ]\leftarrow[X]$\\
    00100001 & 0x21       &  STOR M(X)        & $[X]\leftarrow[AC]$\\
    00000001 & 0x01       &  LOAD M(X)        & $[AC]\leftarrow[X]$\\
    00000010 & 0x02       &  LOAD -M(X)       & $[AC]\leftarrow-[X]$\\
    00000011 & 0x03       &  LOAD |M(X)|      & $[AC]\leftarrow|[X]|$\\
    00000100 & 0x04       &  LOAD -|M(X)|     & $[AC]\leftarrow-|[X]|$\\
    00000101 & 0x05       &  ADD M(X)         & $[AC]\leftarrow[AC]+[X]$\\
    00000111 & 0x07       &  ADD |M(X)|       & $[AC]\leftarrow[AC]+|[X]|$\\
    00000110 & 0x06       &  SUB M(X)         & $[AC]\leftarrow[AC]-[X]$\\
    00001000 & 0x08       &  SUB |M(X)|       & $[AC]\leftarrow[AC]-|[X]|$\\
    00001011 & 0x0B       & MUL M(X)          & $[AC]\leftarrow[MQ]\times[X]$\\
    00001100 & 0x0C       & DIV M(X)          & $[MQ]\leftarrow[AC]/[X],\;[AC]\leftarrow[AC]\bmod[X]$\\
    00010100 & 0x14       &  LSH              & $[AC]\leftarrow[AC]\times2$\\
    00010101 & 0x15       &  RSH              & $[AC]\leftarrow[AC]/2$\\
    \hline
    
  \end{tabular}
\end{table}

\newpage

\subsection{Ejercicio}
\label{sec:org5412a9c}

En la máquina IAS, las instrucciones se dividen en dos segmentos:
i) izquierdo, desde el bit 0 a 19, y ii) derecho, desde el bit 20 al 39. Primero
se ejecuta el lado izquierdo y  luego el derecho. El
contador de programa inicia en la posición 300. El set de
instrucciones del computador o ISA\footnote{Instruction Set Architecture: los códigos de programación} está definido en la Tabla
\ref{tab:tablaISA-IAS}. A fin de terminar el programa se agrega la
instrucción \texttt{0x99} que lleva al computador al estado de HALT. Conteste las
siguientes preguntas considerando la Tabla \ref{tab:orgc59172f} como el
mapa de Memoria de la máquina y obtenga el programa en lenguaje ensamblador y en
RTL. 

\begin{table}[htbp]
\caption{\label{tab:orgc59172f}Mapa de Memoria}
\centering
\begin{tabular}{rrrrr}
\hline
Dirección & \(Opcode_1\) & \(X_1\) & \(Opcode_2\) & \(X_2\)\\[0pt]
\hline
0x300 & 01 & 940 & 06 & 941\\[0pt]
0x301 & 21 & 940 & 14 & 000\\[0pt]
0x302 & 21 & 941 & 99 & 000\\[0pt]
\ldots{} & \ldots{} & \ldots{} & \ldots{} & \ldots{}\\[0pt]
0x940 & 00 & 000 & 00 & 005\\[0pt]
0x941 & 00 & 000 & 00 & 002\\[0pt]
\hline
\end{tabular}
\end{table}


\subsubsection{Preguntas}
\label{sec:org2cb6667}
\begin{enumerate}
\item ¿Qué resultado se tiene al final en el registro del acumulador?
\end{enumerate}
El registro del acumulador [AC], tiene como valor final 0x0000000006
\begin{enumerate}
\item ¿Se sobrescribe algún registro como resultado de la ejecución? Si
verdadero, indique qué registro y con qué valor.
\end{enumerate}
Verdadero, se sobrescriben tres registros: [AC], [940],
   [941]. Quedando con valores finales de:
   [AC] = 0000000006
   [940] = 0000000003
   [941] = 0000000006

\begin{enumerate}
\item ¿Cuál es el valor del bit del Carry?¿Qué operación modifica el bit?
\end{enumerate}
El valor del bit del carry es de 1 y es modificado en la operación
0x06 correspondiente a [AC] <- [AC] - [X], que a su vez corresponde al
[PC] = 0x300 (derecha). En especifico la operación que modifica el
carry es [AC] <- [AC]-[941], [AC] <- 0000000005 - 0000000002

\subsubsection{Código Assembler:}
\label{sec:org59ef1b2}
Escriba el código Assembler del programa que ejecuta el computador IAS.

\begin{verbatim}
LOAD M(940)
SUB M(941)
STOR M(940)
LSH
STOR M(941)
HALT
\end{verbatim}

\subsubsection{Código RTL}
\label{sec:org5312f0c}
Escriba en notación de transferencia de registros el programa que
ejecuta el computador IAS

RTL

\begin{verse}
\([PC] \leftarrow 0x300\)\\[0pt]
\([IR] \leftarrow [PC]\)\\[0pt]
\([IR] \leftarrow 01\ 940\)\\[0pt]
\([AC] \leftarrow [940]\)\\[0pt]
\([AC] \leftarrow 000000005\)\\[0pt]
\([PC] \leftarrow 0x300\ (der)\)\\[0pt]
\([IR] \leftarrow [PC]\)\\[0pt]
\([IR] \leftarrow 06\ 941\)\\[0pt]
\([AC] \leftarrow [AC] - [941]\)\\[0pt]
\([AC] \leftarrow 000000005 - 000000002\)\\[0pt]
\([AC] \leftarrow 000000003\)\\[0pt]
\(Carry = 1\)\\[0pt]
\([PC] \leftarrow 0x301\)\\[0pt]
\([IR] \leftarrow [PC]\)\\[0pt]
\([IR] \leftarrow 21\ 940\)\\[0pt]
\([940] \leftarrow [AC]\)\\[0pt]
\([940] \leftarrow 000000003\)\\[0pt]
\([PC] \leftarrow 0x301\ (der)\)\\[0pt]
\([IR] \leftarrow [PC]\)\\[0pt]
\([IR] \leftarrow 14\ 000\)\\[0pt]
\([AC] \leftarrow [AC] \times 2\)\\[0pt]
\([AC] \leftarrow 000000006\)\\[0pt]
\([PC] \leftarrow 0x302\)\\[0pt]
\([IR] \leftarrow [PC]\)\\[0pt]
\([IR] \leftarrow 21\ 941\)\\[0pt]
\([941] \leftarrow [AC]\)\\[0pt]
\([941] \leftarrow 000000006\)\\[0pt]
\([PC] \leftarrow 0x302\ (der)\)\\[0pt]
\([IR] \leftarrow [PC]\)\\[0pt]
\([IR] \leftarrow 99\ 000\)\\[0pt]
\(Fin \to Estado\ Halt\)\\[0pt]
\end{verse}



\printbibliography
\end{document}
